% Chapter 3: Methodology
% Target: 20-25 pages (5,000-6,000 words)

\section{Problem Formulation}

\subsection{Trading as a Decision Problem}

We formulate the trading problem as follows: Given historical data $D = \{(x_t, y_t)\}_{t=1}^{T}$, where $x_t$ represents market features and $y_t$ represents future returns, the goal is to learn a policy $\pi$ that maximizes risk-adjusted returns.

\subsection{Signal Generation}

We consider two paradigms for signal generation:

\begin{description}
    \item[ML Paradigm] Learn a mapping $f: \mathcal{X} \rightarrow [0, 1]$ from features to signal strength.
    \item[LLM Paradigm] Generate signals through natural language reasoning about market conditions.
\end{description}

\section{Dual-Track Architecture}

\subsection{System Overview}

The Dual-Track framework consists of three main components:

\begin{enumerate}
    \item \textbf{ML Track}: Traditional machine learning pipeline for signal generation
    \item \textbf{LLM Track}: Large language model pipeline with chain-of-thought reasoning
    \item \textbf{Fusion Engine}: Dynamic signal combination based on market state
\end{enumerate}

\subsection{Design Principles}

The framework adheres to several key principles:

\textbf{No Look-ahead Bias}: All features and signals are computed using only information available at time $t$.

\textbf{Fair Comparison}: Both tracks receive identical inputs and are evaluated using the same metrics.

\textbf{Reproducibility}: All experiments are fully reproducible with fixed random seeds.

\section{ML Track Implementation}

\subsection{Feature Engineering}

We compute 50+ technical indicators organized into categories:

\begin{table}[h]
\centering
\caption{Feature Categories and Examples}
\begin{tabular}{lll}
\toprule
Category & Indicators & Count \\
\midrule
Momentum & RSI, MACD, ROC & 15 \\
Volatility & ATR, Bollinger Bands & 10 \\
Volume & OBV, VWAP & 8 \\
Trend & MA, EMA, ADX & 12 \\
Others & & 10 \\
\bottomrule
\end{tabular}
\end{table}

\subsection{Baseline Models}

We implement three baseline models:

\subsubsection{Logistic Regression}

A linear model serving as a simple baseline:
\begin{equation}
P(y_t=1|x_t) = \sigma(w^T x_t + b)
\end{equation}

\subsubsection{LightGBM}

Gradient boosting decision trees \cite{ke2017lightgbm} with optimized hyperparameters...

\subsubsection{LSTM}

Long Short-Term Memory networks \cite{hochreiter1997lstm} for sequence modeling...

\section{LLM Track Implementation}

\subsection{Prompt Engineering}

We design chain-of-thought prompts following a structured framework:

\begin{enumerate}
    \item Market context description
    \item Technical indicator summary
    \item News sentiment analysis
    \item Risk assessment
    \item Trading decision with reasoning
\end{enumerate}

\subsection{Model Selection}

We evaluate multiple LLM backends:

\begin{itemize}
    \item DeepSeek V3.2 (cloud)
    \item DeepSeek V3.2 Reasoning (cloud)
    \item DeepSeek R1 14B (local)
    \item DeepSeek R1 8B (local)
\end{itemize}

\subsection{Output Parsing}

LLM outputs are parsed using robust JSON extraction:

\begin{itemize}
    \item Signal strength extraction
    \item Confidence estimation
    \item Reasoning chain capture
\end{itemize}

\section{Fusion Engine}

\subsection{Market State Classification}

We classify market states based on volatility:

\begin{table}[h]
\centering
\caption{Market State Classification}
\begin{tabular}{lccc}
\toprule
State & Volatility & ML Weight & LLM Weight \\
\midrule
Normal & $< 2\%$ & 70\% & 30\% \\
High Volatility & $2\% - 5\%$ & 50\% & 50\% \\
Black Swan & $> 5\%$ & 0\% & 100\% \\
\bottomrule
\end{tabular}
\end{table}

\subsection{Signal Combination}

The fusion engine combines signals using a weighted approach:

\begin{equation}
s_t = w_{ML} \cdot s_t^{ML} + w_{LLM} \cdot s_t^{LLM}
\end{equation}

\section{Reproducibility Considerations}

\subsection{Random Seeds}

All experiments use fixed random seeds for reproducibility...

\subsection{Model Persistence}

Trained models are saved with metadata...